\documentclass[UTF8,fontset=none,11pt,a4paper]{ctexart}
\usepackage{ipara-handbook}
\handbooksetup{NET-B05 BDP × Window}{408 · Network Internal Bridge}{Rate · RTT · In-Flight · Utilization}
\begin{document}
\handbooktitle{BDP × Window}{物理管道容量怎样生成窗口需求}{408 · NET-B05}

\begin{corebox}{Mother Interface}
\[
\boxed{\text{Path rate/time}\to\text{required in-flight data}\to\text{window bound}\to\text{utilization}}
\]
BDP 是某段传播期间可注入的数据量；可靠窗口是在反馈到达前允许保持的未确认数据量。两者通过“反馈管道需要多少数据才能不空”连接。
\end{corebox}

\begin{methodbox}{Owners 与边界}
NET01 Owns $R,T_{prop},RTT$ 与 delay decomposition；NET03 Owns window、ACK 与流水线。本桥 Owns 量纲转换和最小窗口条件。\code{rwnd/cwnd} 如何限制实际窗口交给 NET-B04；队列拥塞交给 NET07。
\end{methodbox}

\section{先统一时间口径}

若只问“单向链路上同时传播的数据”，则 $BDP=R\,T_{prop}$。若问发送方从开始发送一个 frame 到其 ACK 返回前的停等周期，就要包含题设中的发送、往返传播、ACK 发送及其他处理时间：
\[
T_{cycle}=T_{tx,data}+2T_{prop}+T_{tx,ack}+T_{other}.
\]
窗口要让发送器在整个反馈周期内持续忙碌，所需数据量近似为 $R\,T_{cycle}$；如果题目把“窗口”只计未确认 data frames，还应按 frame 长度离散取整。

\section{从停等到流水线}

设 data frame 长 $L$ bit，发送时间 $T_f=L/R$，忽略 ACK 发送与处理。停等利用率为
\[
U_{SW}=\frac{T_f}{T_f+2T_{prop}}.
\]
窗口允许最多 $W$ 个 frame 在途时，理想化发送端利用率为
\[
U=\min\left\{1,\frac{W T_f}{T_f+2T_{prop}}\right\}.
\]
要达到满利用率：
\[
W\ge\left\lceil\frac{T_f+2T_{prop}}{T_f}\right\rceil
=\left\lceil1+2\frac{T_{prop}}{T_f}\right\rceil.
\]
式中的 $1$ 对应首帧本身的发送时间；若题目另定义 RTT 从“首帧发送完”起算，公式外观会变，但事件时间线必须等价。

\section{BDP 只是必要容量，不是实际吞吐保证}

窗口达到反馈管道需求，只消除了“等 ACK 导致发送器空闲”这一瓶颈。实际吞吐还可能受 source rate、receiver \code{rwnd}、congestion \code{cwnd}、loss/retransmission、最慢链路和协议开销限制：
\[
Throughput\le\min\left(R_{bottleneck},\frac{W_{effective}}{RTT}\right).
\]
这是一条上界接口，不是所有题目都能直接取等。RTT 增大时，同一窗口能支撑的吞吐下降；链路速率增大但窗口不变时，也可能出现“带宽升级、吞吐不升”。

\section{数值母例}

链路 $R=100$ Mbps，单向传播 $T_{prop}=20$ ms，data frame 为 $10{,}000$ bit，忽略 ACK：$T_f=0.1$ ms，反馈周期约 $40.1$ ms，需要
\[
W_{min}=\left\lceil\frac{40.1}{0.1}\right\rceil=401\text{ frames}.
\]
对应约 $4.01$ Mbit 未确认数据。只计算单向 $R T_{prop}=2$ Mbit 会低估可靠发送方等待 ACK 的在途需求，因为反馈要走完往返路径。

\section{调用协议与 First Divergence}

\begin{enumerate}
\item 画“开始发首帧—发完—到达—ACK 返回”时间线；
\item 标明 RTT/传播时延的题设定义，列入 ACK 与处理开销；
\item 把窗口统一成 bit、byte 或 frame，做量纲检查并向上取整；
\item 算满利用率所需窗口，再与题设实际有效窗口比较；
\item 若未达理论吞吐，继续检查 NET-B04 与 bottleneck，而非强行取等。
\end{enumerate}

\begin{warnbox}{Anti-Bridge}
$BDP$ 当成一种时延，是量纲错误；用单向 BDP 直接回答 ACK 窗口，是漏了反馈返回；窗口大于 BDP 就断言不会拥塞，是把物理填管需求变成网络容量许可；忘记 frame 离散取整，会少一个窗口槽。
\end{warnbox}

\begin{corebox}{Compression}
窗口覆盖的是单向传播还是完整反馈周期？速率与时间的口径是什么？需要多少在途 bit/frame？实际窗口又被谁进一步收紧？
\end{corebox}
\end{document}
